\section{Payments}
\label{sec:payments}

All payments are in bitcoin. There is no swap service, mixer, ecash mint, or second asset in the design or the code. Direct payment goes from the buyer to the vendor with no intermediary; moderated payment goes through the escrow of \S\ref{sec:escrow}.

\subsection{One payment per seller}

A multi-seller checkout produces one order per seller, and each order is paid by its own transaction \IT{} (PR~\#2). The pay-from-node path builds a transaction with exactly one recipient output and, if needed, one change output back to the buyer's own wallet. Two sellers are never paid in one transaction, because a transaction with outputs to two vendors would tell anyone watching the chain that the same person bought from both.

Separate transactions are necessary for unlinkability, not sufficient. If the wallet funds the second payment with the change from the first, the common-input-ownership and change heuristics~\citep{meiklejohn2013fistful} link the two payments. At \paycommit{} coin selection skips unconfirmed coins received from others but deliberately keeps the wallet's own unconfirmed change spendable, so paying two sellers in quick succession can chain the second payment onto the first payment's change. The wallet does not currently avoid this; choosing inputs so that one checkout's payments do not share ancestry is future work \FW.

\subsection{On-chain direct payment}

In PR~\#2, each direct order receives a fresh receive address from the vendor's BDK descriptor wallet (BIP 86 Taproot) and a BIP 21~\citep{bip21} URI for the exact amount \IT. The node's chain watcher queries an Esplora server for transactions paying each watched address and marks an order as payment seen and then paid once the required confirmations (one on testnet4, three by default on mainnet) are reached. The node can also pay an order from its own wallet. The default network is testnet4~\citep{bip94}. Mainnet is selected only by the exact flag \code{--network mainnet}, and a mainnet payment from the node additionally requires \code{"confirm": true} in the request.

Because the demonstration node holds both the buyer wallet and the vendor wallets, a single node can both pay and observe an order \ID. In a deployment the buyer and vendor nodes are different machines, and each watches only its own addresses.

\paragraph{The Esplora server.} Every chain query names a specific address to the Esplora server~\citep{esplora}. The defaults are \code{mempool.space} on mainnet and, on testnet4, \code{mempool.space} with \code{blockstream.info} as an ordered fallback; the node checks that each server's genesis block matches its network before using it. Whichever server answers learns the full set of addresses the node watches, their timing, and the IP address of the node, which together link a user's orders to each other and to a network location; once the node falls back, the second operator learns the same. On mainnet the node prints a warning when the server is public. Routing chain queries over Tor \SP{} would hide the IP address but not the address set; running one's own Esplora or Bitcoin Core node, or using compact block filters, would hide both. None of these is implemented.

\subsection{Lightning: BOLT 12}

A store record carries a BOLT 12 offer~\citep{bolt12}, and direct checkout returns an offer for the order's exact amount \status{encoding only}. The offer is a valid \code{lno1} string built with LDK's offer builder, signed with the store's signing public key, and carrying no blinded path. No Lightning node in the system answers invoice requests for it, so it cannot be paid today.

Two privacy points apply once it can. An offer without a blinded path names the issuer's node key, so the offer links the store to a Lightning node. Using the store signing key as the offer's issuer key, as the current code does, also links the catalog identity to that Lightning identity; a separate key is the intended design \SP. On a payment, every routing hop sees the amount it forwards and its neighbors; the recipient's position is hidden from the payer only to the extent that a blinded path of more than one hop is used~\citep{kappos2021lightning}.

\subsection{Silent payments: BIP 352}

The store record also carries a BIP 352 silent-payment code~\citep{bip352} \status{encoding only}. A silent-payment code lets a payer derive a fresh output for each payment from a static public code, so a single published code does not lead to address reuse. The code is constructed from the scan and spend public keys and parsed with checksum validation. The node does not scan the chain for payments to the code. A buyer who pays the code today creates an output the vendor's node will not detect until scanning exists; the funds are recoverable with the scan and spend keys but the order would not be marked paid. The intended design keeps the scan key, which is watch-only, on the vendor's node so that scanning does not hand the key to a public indexer \SP.

\subsection{Payjoin: BIP 77}

Payjoin v2~\citep{bip77} lets a receiver add an input to the buyer's transaction through an asynchronous directory reached over Oblivious HTTP~\citep{rfc9458}, which breaks the assumption that all inputs belong to the payer. The \code{pay} crate wraps the \code{payjoin} library's v2 sender and receiver sessions and the step where the vendor contributes an input \status{library only}. No checkout starts a Payjoin session, and the payment terms carry no \code{pj=} endpoint. When Payjoin is used, the vendor learns which coins the buyer spent, and the buyer learns one of the vendor's coins.

\subsection{What the chain shows}

For an on-chain direct payment, a chain observer sees an ordinary transaction with one or two outputs. The amount is the listing price times quantity plus shipping. Prices are public, so an amount that matches a distinctive price can identify the listing that was bought, even though nothing else in the transaction refers to OpenBazaar. Round prices reduce this and do not remove it. For escrow, see \S\ref{sec:escrow-privacy}.
